\documentclass[11pt]{article}
\usepackage{fontspec}
\usepackage{xeCJK}
\setCJKmainfont[ItalicFont=STSongti-SC-Bold]{Songti SC}
\usepackage{xurl}
\usepackage[a4paper,margin=25mm]{geometry}
\usepackage{amsmath,amssymb,amsthm,mathtools}
\usepackage{xcolor}
\usepackage{booktabs,array,enumitem,fancyhdr,hyperref}
\hypersetup{colorlinks=true,linkcolor=blue!50!black,urlcolor=blue!50!black}
\setlength{\emergencystretch}{2em}
\setlength{\parskip}{0.3em}
\linespread{1.16}
\allowdisplaybreaks
\newtheorem{theorem}{定理}[section]
\newtheorem{proposition}[theorem]{命题}
\newtheorem{lemma}[theorem]{引理}
\theoremstyle{definition}
\newtheorem{definition}[theorem]{定义}
\theoremstyle{remark}
\newtheorem{remark}[theorem]{说明}
\newcommand{\dd}{\,\mathrm d}
\newcommand{\ii}{\mathrm i}
\newcommand{\Rea}{\operatorname{Re}}
\newcommand{\Ima}{\operatorname{Im}}
\newcommand{\Li}{\operatorname{Li}}
\newcommand{\supp}{\operatorname{supp}}
\pagestyle{fancy}
\fancyhf{}
\fancyhead[L]{\small 学习笔记}
\fancyhead[R]{\small 初稿}
\fancyfoot[C]{\thepage}
\setlength{\headheight}{15pt}
\author{}
\date{2026年10月22日初稿}
\makeatletter
\renewcommand{\maketitle}{%
  \thispagestyle{plain}%
  \begin{center}
    {\LARGE\bfseries \@title\par}
    \vspace{0.6em}
    {\small \@date\par}
  \end{center}
  \vspace{0.6em}}
\makeatother
\newcommand{\draftnotice}{\noindent{\small
本文为供逐篇审阅的学习笔记初稿，不是新增研究成果。
文中区分直接证明与引用结果；正文不依赖本文。
本轮已作公式、条件及排版初步核查，尚待作者逐句审定。}\par\medskip}

\XeTeXlinebreaklocale "zh"
\renewcommand{\proofname}{证明}
\date{2026年10月24日初稿}
\title{数值证据在黎曼猜想中的逻辑地位}
\begin{document}
\maketitle

\noindent{\small 本篇讨论零点数值验证能确立什么、不能确立什么。凡标为命题者给出证明；
引用的定理与计算记录注明出处；标为“启发性”的论证不作为证明使用。}\par\medskip

\section{问题}
Platt--Trudgian（2021）以严格认证的计算证明：高度不超过 $3\cdot10^{12}$ 的全部非平凡零点
（共 $12\,363\,153\,437\,138$ 个）都位于临界线上。
这一事实常被表述为支持 RH 的“压倒性证据”。本篇要回答三个问题：
\begin{enumerate}
 \item 这类计算在逻辑上证明了什么；
 \item 它们能否转化为关于全体零点或素数分布的严格结论；
 \item 在何种尺度上，这些数据是“大量”的，在何种尺度上则微不足道。
\end{enumerate}

\section{有限验证的逻辑形式}
记 $\rho=\beta+\ii\gamma$ 为非平凡零点。对 $T>0$，令
\[
 \mathrm{RH}(T):\quad \text{凡 } 0<\gamma\le T \text{ 的非平凡零点均满足 } \beta=\tfrac12 .
\]
由零点关于实轴对称，$\mathrm{RH}$ 等价于 $\forall T\,\mathrm{RH}(T)$。

认证一个给定高度时，可以先由 Hardy 函数 $Z(t)$ 的严格变号得到临界线上零点个数的下界，
再用 Turing 方法核对临界带中的按重数总计数（学习笔记第 16 篇）。
若两数相等，就证明了这一高度内的 RH，且所计零点均为单零点。
这不是对任意 $T$ 都保证成功终止的判定算法；例如多重零点未必由简单的变号检测取尽。
Platt--Trudgian 的计算完成了所需核验，因而 $\mathrm{RH}(3\cdot10^{12})$ 是定理，而非经验观察。

然而，仅凭有限多个 $\mathrm{RH}(T_j)$、而没有控制其余高度的一般论证，不能推出 $\forall T\,\mathrm{RH}(T)$；
反之，若 RH 不成立，离线零点的存在可以由有限计算证实。这一不对称的严格表述见学习笔记第 13 篇。

\section{部分验证转化为严格定理}
有限验证虽不能证明 RH，却能作为严格证明的输入，产生关于全体零点或素数的无条件结论。

\subsection{对显式公式的作用}
回忆显式公式中零点项 $\sum_\rho x^\rho/\rho$。
\Needspace{9\baselineskip}
\begin{proposition}\label{prop:lowzeros}
设 $\mathrm{RH}(T)$ 成立。则对任意 $x>1$，
\[
 \Bigl|\sum_{|\gamma|\le T}\frac{x^\rho}{\rho}\Bigr|\le
 x^{1/2}\sum_{|\gamma|\le T}\frac1{|\gamma|}.
\]
\end{proposition}
\begin{proof}
由 $\mathrm{RH}(T)$，所列零点满足 $|x^\rho|=x^{1/2}$，且 $|\rho|\ge|\gamma|$。逐项取模即得。
\end{proof}
由 $N(T)$ 的渐近式，$\sum_{0<\gamma\le T}\gamma^{-1}=\tfrac1{4\pi}\log^2T+O(\log T)$。
因此低处零点对 $\psi(x)-x$ 的贡献得到 RH 级别的控制；
高处零点的贡献则由零点密度估计与无零区域另行处理。
这正是用部分验证改进 $\pi(x)$、$\psi(x)$ 显式误差界的基本思路
（例如 Büthe 2018；D.\ R.\ Johnston 2022 利用 $\mathrm{RH}(3\cdot10^{12})$ 改进了 $\pi(x)$ 的显式界）。

\subsection{对 de Bruijn--Newman 常数的作用}
记 $\Lambda$ 为 de Bruijn--Newman 常数，RH 等价于 $\Lambda\le0$。
Rodgers 与 Tao（2020）证明 $\Lambda\ge0$。Polymath15 项目证明了如下类型的结论：
若 $\mathrm{RH}(T)$ 对足够大的 $T$ 成立，并辅以有限范围内的计算，则可得 $\Lambda$ 的上界。
由此得到 $\Lambda\le0.22$（2019），Platt 与 Trudgian 结合 $\mathrm{RH}(3\cdot10^{12})$ 改进为 $\Lambda\le0.2$（2021）。
这是数值验证直接转化为关于全体零点的定理的一个例子：
$\Lambda\le0.2$ 是一个刻画所有零点的命题。

\section{小范围规律失效的先例}
在解析数论中，有限范围内成立的规律曾多次在更大范围内失效：
\begin{itemize}
 \item \textbf{$\pi(x)<\operatorname{li}(x)$}。对 $2\le x\le10^{19}$ 成立（Büthe 2018），
 但 Littlewood（1914）证明 $\pi(x)-\operatorname{li}(x)$ 无穷多次变号；
 第一次变号的位置已知小于 $1.4\times10^{316}$（Saouter--Demichel 2010 等）。
 \item \textbf{Pólya 猜想} $L(x)=\sum_{n\le x}\lambda(n)\le0$（$x\ge2$，$\lambda$ 为 Liouville 函数）。
 Haselgrove（1958）证明其不成立；最小反例为 $x=906\,150\,257$（Tanaka 1980）。
 \item \textbf{Mertens 猜想} $|M(x)|<\sqrt x$（$x>1$）。Odlyzko 与 te Riele（1985）证明其不成立，
 但没有给出具体反例；Rozmarynowycz--Kim（2023）给出的反例上界为 $\exp(1.017\times10^{29})$。
\end{itemize}
三者都通过显式公式与零点相联系，且失效位置都在早期数据所及范围之外。
以对数变量计，所列失效位置或其上界的 $\log x$ 依次约为 $728$、$20.6$、$1.02\times10^{29}$。
后两例的差距尤其说明，失效位置可以跨越极不相同的尺度。

\subsection{多项相位对齐：存在界与长期占比}
在对数变量 $u=\log x$ 中，零点项是 $e^{\ii\gamma u}$ 型的振荡（学习笔记第 7 篇）。
某些大幅振荡的构造需要同时控制多项的相位。
先看不带系数、目标相位全为 $1$ 的模型：下述命题给出一次近似对齐的存在上界。
\Needspace{12\baselineskip}
\begin{proposition}[Dirichlet 联立逼近]\label{prop:dirichlet}
设 $\gamma_1,\ldots,\gamma_K$ 为实数，$Q\ge1$ 为整数。
则存在整数 $1\le q\le Q^K$，使 $u=2\pi q$ 满足
\[
 |e^{\ii\gamma_ku}-1|\le\frac{2\pi}{Q},\qquad k=1,\ldots,K .
\]
\end{proposition}
\begin{proof}
考虑 $Q^K+1$ 个点 $(\{j\gamma_1\},\ldots,\{j\gamma_K\})\in[0,1)^K$，$j=0,1,\ldots,Q^K$。
把 $[0,1)^K$ 分成 $Q^K$ 个边长 $1/Q$ 的小方体，由抽屉原理有 $j<j'$ 落入同一方体。
令 $q=j'-j$，则 $q\gamma_k$ 与最近整数之差小于 $1/Q$。
由 $|e^{2\pi\ii\theta}-1|\le2\pi\|\theta\|$ 即得。
\end{proof}
命题只给出上界：$K$ 个相位对齐到精度 $2\pi/Q$，不迟于 $u\le2\pi Q^K$。
若再假设 $\gamma_1,\ldots,\gamma_K$ 在 $\mathbb Q$ 上线性无关，
连续参数的 Kronecker--Weyl 等分布定理给出另一类信息。固定 $0<\delta<2$，令
\[
 E_U=\{u\in[0,U]: |e^{\ii\gamma_ku}-1|\le\delta\ (1\le k\le K)\}.
\]
则 Lebesgue 测度满足
\[
 \frac{|E_U|}{U}\longrightarrow
 \left(\frac{2\arcsin(\delta/2)}\pi\right)^K.
\]
因为单位圆上到 $1$ 的距离不超过 $\delta$ 的弧长占比为
$2\arcsin(\delta/2)/\pi$，小 $\delta$ 时才近似为 $\delta/\pi$。
这说明固定精度下的长期占比可以随 $K$ 很快减小，
但不提供首次进入该集合的时间界，也不保证有限窗口已经表现出这一比例。
尤其在 $u=0$ 附近，任意有限组相位本来就同时接近 $1$。

零点虚部的线性无关性本身尚是猜想；实际显式公式还含有复系数与未截取的尾项，
促成特定变号的目标相位也未必全为 $1$。
因此不能据此断言 $x\le10^{19}$ 时只有少数零点能够协同作用，
或据长期占比预测 RH 反例的高度。这里的模型仅解释多相位条件为何需要单独分析。

\section{零点验证的范围放在合适的尺度上}
零点数据的情形与上节不同：验证的对象是零点本身，而不是由零点决定的算术函数。
但零点的某些统计量同样随高度极慢地变化。

记 $S(t)=\frac1\pi\arg\zeta(\tfrac12+\ii t)$（沿通常约定的路径连续延拓）。
它度量 $N(t)$ 与其光滑主项之差，是累计计数偏差。
长度 $h$ 的短区间内零点数还涉及 $S(t+h)-S(t)$，不能仅由单点的 $|S(t)|$ 描述。
\begin{theorem}[Selberg 中心极限定理，引用]
当 $T\to\infty$，对 $t$ 在 $[T,2T]$ 上均匀分布，
$S(t)\big/\sqrt{\tfrac1{2\pi^2}\log\log T}$ 依分布收敛于标准正态分布。
\end{theorem}
在 $T=3\cdot10^{12}$ 处，$\log T\approx28.73$，$\log\log T\approx3.358$，
相应的标准差约为 $0.41$。已知无条件地有 $S(t)=O(\log t)$；
在 RH 下有 $S(t)=O(\log t/\log\log t)$，并且 $S(t)$ 无界
（无条件地，Selberg 证明了 $S(t)$ 取到 $\gg(\log t)^{1/3}(\log\log t)^{-4/3}$ 量级的值）。

因此，从零点个数看，$1.2\times10^{13}$ 个零点是很大的样本；
从上述特定中心极限定理所用的变量 $\log\log T$ 看，该连续验证高度仅达到约 $3.4$。
若存在离开临界线的零点，它们与 $|S(t)|$ 的大值并无已知的必然联系；
上述比较不说明有限计算只探查了典型行为，也不提供极端行为出现高度的预测；
它只是提醒我们，不同统计命题使用的自然尺度不同，零点个数不能替代对适用范围的分析。

\Needspace{15\baselineskip}
\section{结论}
\begin{enumerate}
 \item 经认证的零点计算证明的是 $\mathrm{RH}(T)$ 这一有限命题。它是定理，而非经验。
 \item 它不能推出 RH；但 RH 的否定若成立，原则上可以由有限计算证实。
 因此数值验证排除了“较低处的反例”，这本身是确定的信息。
 \item 部分验证可以作为严格证明的输入，产生关于全体零点（如 $\Lambda\le0.2$）
 与素数分布（如显式误差界）的无条件定理。
 \item 评价大规模数据须同时交代认证对象与所讨论的尺度。
 算术函数变号、相位条件和 $S(t)$ 的统计是不同问题；
 它们可以帮助辨认推断的边界，却不能用来估算离线零点出现的概率或位置。
\end{enumerate}

\section*{参考}
\begin{itemize}
 \item D.\ Platt, T.\ Trudgian, The Riemann hypothesis is true up to $3\cdot10^{12}$,
 \emph{Bull.\ London Math.\ Soc.} 53 (2021). \url{https://arxiv.org/abs/2004.09765}
 \item B.\ Rodgers, T.\ Tao, The de Bruijn--Newman constant is non-negative, \emph{Forum Math.\ Pi} 8 (2020).
 \item D.\ H.\ J.\ Polymath, Effective approximation of heat flow evolution of the Riemann $\xi$ function, and a new upper bound for the de Bruijn--Newman constant, \emph{Res.\ Math.\ Sci.} 6 (2019).
 \item J.\ Büthe, An analytic method for bounding $\psi(x)$, \emph{Math.\ Comp.} 87 (2018).
 \item D.\ R.\ Johnston, Improving bounds on prime counting functions by partial verification of the Riemann hypothesis, 2022. \url{https://arxiv.org/abs/2109.02249}
 \item A.\ M.\ Odlyzko, H.\ J.\ J.\ te Riele, Disproof of the Mertens conjecture, \emph{J.\ reine angew.\ Math.} 357 (1985).
 \item J.\ Rozmarynowycz, S.\ Kim, A new upper bound on the smallest counterexample to the Mertens conjecture, 2023. \url{https://arxiv.org/abs/2305.00345}
 \item A.\ Selberg, Contributions to the theory of the Riemann zeta-function, \emph{Arch.\ Math.\ Naturvid.} 48 (1946).
 \item E.\ C.\ Titchmarsh, \emph{The Theory of the Riemann Zeta-Function}, 2nd ed., 1986, Ch.~IX.
\end{itemize}
\end{document}
